\subsection{封国不是一排听令的据点}

\begin{event}{周初至成康时期}{重组东方封国网络}{总体过程可靠；各国始封的年月、受封者与边界常有争议}{成周以东及各地政治节点}\eventanchor{WZ-EAST-ENFEOFFMENT}{西周·重组东方封国网络}
东征的军队撤去之后，东方留下的不是一张可以交接的地图，而是许多仍在运作的宗族、旧商贵族、聚落与交通通道。周王室必须决定：哪些人可被托付，哪些旧有关系应被拆开，哪些地方要保留足以集结粮食、兵员和祭祀的政治节点。宗亲、功臣和被重新纳入的旧商贵族因此陆续进入不同位置；后世所说的鲁、齐、燕、卫、宋等封国，正是在这段漫长而不齐整的重组中获得了可继承的政治名分。

传世叙事喜欢把这些安排写成一连串清楚的授封：谁在某年受土，谁在何处立国，疆界又到哪里。然而，能较早看见的材料并不提供这样的总册。\sourcebook{尚书·康诰}、\sourcebook{酒诰}与\sourcebook{梓材}保存了对康叔及殷民的告诫和治理语言，却不是附有地籍的封授文书；\sourcebook{史记}各世家把鲁、齐、燕、卫、宋的祖源串成较完整的谱系，成书却远在事件之后。鲁侯簋、燕侯相关金文和琉璃河等地的早周遗存，能让贵族礼器、墓地和地方经营从纸面走到地下，却仍不能替每一国签出确切的始封日期、第一道边界或一份移民名册。

因此，较稳妥的说法不是“周初一次分完了东方”，而是王室在\eventref{WZ-1040-DONGZHENG}{西周·三监东征}之后，以册命、婚姻、驻守、祭祀和资源调动反复确认若干地方节点。对受命家族而言，这带来了可传给后代的地位；对当地原有共同体而言，则意味着须在新来的政治中心、旧有生计与服役要求之间重新安排关系。它的直接后果，是东方不再只靠一支远征军维持联系；它的中期代价，是王室必须不断用赏赐、裁断和亲临，使这些各有根基的节点仍愿接续王命。
\end{event}

成周的建设必须和东方封国的布置一起看。后来常用“分封制”概括周初安排；这一概念有助于说明王室把土地、名分和防务交给不同政治集团，却容易让人误以为当时已经有一张整齐的行政区划图。周王直接掌握的王畿、宗室支系的封地、功臣和盟友的政治基础，以及原有地方共同体之间，实际存在不同层次的关系。各地形成的先后、疆界与王室控制程度也并不相同。

王室提供的不是单纯的一块土地，而是一套可被承认的资格：受命者可以在当地组织宗族、祭祀和武装，并以周王授予的名分来处理对内对外关系。作为交换，王室期待他们在朝见、用兵、礼仪和资源供给上承担义务。这里的“期待”尤其重要。西周没有一套留存至今、可逐条核对的统一地方税法；不同封国在何时以何种物资、兵员或礼仪回应王室，不能用后世官僚制的标准倒推。

这是一种低成本却不廉价的统治方法。王室不必在每一处都长期派驻同等规模的官员和军队，封国的宗族资源能够承担守卫、开垦、调解与扩张；但中心要持续确认其地位，处理彼此冲突，并以赏赐、婚姻和仪式维持可见的等级。王权的力量不只是能否下令，更在于这些地方精英是否仍把王室的命令视为值得兑现的名分。

\begin{causalchain}[从东方据点到关系网络]
\term{结构性原因}是周初中心人力有限、东方政治关系复杂，不能只靠从关中远征来治理。\term{使能条件}是宗族、军功、礼器与祭祀能把地方精英编进共同语言。成周提供了较近的会集点，封国则提供日常执行者；二者结合，才形成\eventref{WZ-EAST-CHENGZHOU}{西周·营建成周}所代表的东方网络。\term{反馈回路}也在这里产生：王室每次对功劳的承认，都增强了受命者的资源和声望；资源越集中于地方，王室日后越要依赖协商，而不能只凭命令。
\end{causalchain}

\subsection{从册命到地方能力}

现存青铜器铭文经常记载王命、赏赐、服事和祭祀，说明名分并非抽象口号：它要通过反复的仪式和器物被保存、展示并传给后人。不过，这类材料大多由获得命令或赏赐的贵族制作，保存的是成功进入王室关系网的人如何叙述自己，较少留下未获承认者、普通服役者或地方对抗者的声音。

因此，东方网络不能只用“中央强、地方弱”或反过来的二分法理解。在\eventref{WZ-CHENGKANG-CONSOLIDATION}{西周·成康巩固}所代表的阶段，王室权威能够为封国提供共同的秩序语言，封国也把本地人口、土地和武力接入更广的体系。长远看，同一机制又让地方家族积累连续的祭祀传统、军事组织和资源调动能力。西周末年的危机并非由周初封国单独造成，却会在这些已具实力的地方网络中寻找盟友、兵员和退路。

\begin{debate}[“封建”能解释多少]
“封建”是后世整理周代制度时形成的概念，适合描述授土、建国和宗族政治彼此相连的现象，却不应当替代对具体地区的考察。研究者会结合金文中的册命和赏赐、传世文献的追述、聚落与墓葬的区域差异，讨论王室控制如何随地点和时期改变。没有一件材料能单独列出完整封国清单，更不能据此断言所有封国都以同一种方式服从王室。
\end{debate}

\subsection{几条道路，不是一份整齐的封国名册}

\begin{event}{周初至成王时期}{伯禽受封鲁地并经营曲阜据点}{受封关系可信；初期疆域及经营细节不定}{曲阜及其联通的东部聚落}\eventanchor{WZ-EARLY-LU-FOUNDATION}{西周·伯禽经营鲁地}
周公的家族若要在东部留下可继承的据点，不能只依赖宗庙中的祖先名分。后世叙事把伯禽受封鲁地写得连贯，曲阜周公庙一带的遗存则提示这里确有值得长期经营的中心。二者合在一起，支持鲁成为连接周人与地方人群的东部支点；它并不使我们得到一份周初曲阜的户口册或疆界图。

王室从鲁得到的，是一个能在东部组织祭祀、武力和地方关系的宗亲节点；伯禽及其后代得到的，则是必须在当地维持的政治资格。这个交换的成本落在周人和当地共同体的重新接合上。谁迁入、谁保有土地、哪些礼俗在哪个阶段发生改变，现有材料不能逐项回答，故不应把“受封”缩写为一纸授土即完成的事实。
\end{event}

围绕伯禽“就国”的故事，\eventanchor{WZ-1030S-BOQIN-LU-SETTLEMENT}{西周·伯禽就国的叙事}又把礼俗整合写成一段可供后世反复讲述的鲁国起源。\sourcebook{史记·鲁周公世家}与\sourcebook{汉书·地理志}保存的是较晚的叙事和地理记忆，不足以坐实某次会面、某项礼制或确切年代。它仍有历史价值：它表明后来的鲁人和史家把封国的持久性理解为周人礼仪与本地人群之间的调适，而非仅仅一次军事占领。

\begin{event}{周初至成王时期}{姜尚家族受封营丘一带并成为山东通道的节点}{谱系叙事可用；营丘的具体定位仍有争议}{营丘一带及山东交通方向}\eventanchor{WZ-EARLY-QI-FOUNDATION}{西周·姜尚家族受封营丘}
齐的早期叙事把姜尚家族的军功、受命和营丘联系起来。这种联系说明周王室并不只用宗亲布置东方：功臣网络也能承担通向山东的关系与资源。临淄一带的早期遗存使地方中心的长期形成不只依赖后世谱系，但遗存不能自动标出“营丘”，更不能由此画定齐的最初边界。

因而，齐的建国故事应被读作军功、资源和名分相互支撑的传统。它使一个家族获得在当地积累人手和物资的资格，也让王室能够借其进入山东通道；至于受封发生的单一年份、所在地是否可精确对应、当地各群体如何参与，仍须让材料的空白保留在叙事中。
\end{event}

\begin{event}{西周早期传统}{唐叔虞受封唐、后称晋的谱系叙事形成}{传统可以追溯；改称和早期发展过程未定}{汾河流域及其交通资源}\eventanchor{WZ-1030S-JIN-TANGSHU}{西周·唐叔虞受封唐的传统}
向东北、北方展开的网络也不只靠燕地。关于唐叔虞受封唐、后来称晋的说法，\sourcebook{史记·晋世家}给出了可读的家族起点；天马—曲村的墓地及晋侯器则使一个持续祭祀的地方精英传统进入考古视野。两类材料可以彼此照亮，却不能把较晚的谱系逐句倒译成西周早期的政治地图。

较可把握的是，汾河流域的土地、交通和地方资源需要由宗亲据点加以联结。唐、晋名称的转换何时发生、早期中心如何移动、哪些墓主可与世系一一相配，都是未决问题。正因为这些问题未决，晋的形成更适合被写作一个谱系叙事与地方经营相互叠合的过程。
\end{event}

\begin{event}{康王时期说}{宜侯夨簋保存王命、赐土与宜侯事务}{器铭可靠；地名性质及其与后世地域的关系有争议}{东南方向的受命关系}\eventanchor{WZ-KANG-YIHOU-ZE-GUI}{西周·宜侯夨簋的东南配置}
宜侯夨簋比许多后出的建国故事更接近当时的受命场合。器铭将王命、赐与和宜侯事务并置，显示边缘据点也要以王室认可来确认承担者。它不必意味着一个现代行政单位在此刻被完整设立，却让东方、东南的政治配置留下了具体的金文证据。

“宜”的地望及其与后世吴地叙事的关系仍有争议。\sourcebook{史记·吴太伯世家}可以提供后来的谱系框架，不能替器铭裁断地点。由此可见的是王命能触及的关系网络；看不见的仍包括当地人群的构成、赐土的实际范围和受命者日后怎样把命令落实为地方权力。
\end{event}

\subsection{墓地保存的时间，比授封故事更长}

\begin{event}{西周早期至晚期}{天马—曲村晋侯墓地显示家族祭祀的延续}{层位和遗存可靠；墓主与传世世系的对应仍有讨论}{晋地的墓地、礼器与地方资源}\eventanchor{WZ-EARLY-JIN-MARQUIS-CEMETERY}{西周·天马—曲村晋侯墓地}
一座跨越多期的墓地，保存的不是王室每一次命令，却能看见家族如何把资格留给后人。天马—曲村墓地的层位和晋侯器表明，地方精英持续以祭祀、礼器和丧葬组织自身；这比单一的始封故事更接近封国能力如何累积的物质过程。王室认可的名分，需要在地方的土地、人力和仪式中一代代兑现。

墓主究竟可对应传世世系中的哪一位晋侯，仍不能处处确定。因此，墓地不能用来校正每一个王年，也不能直接测量晋国的军力。它所支持的较小却坚实的判断是：宗亲政治体拥有长期经营家族记忆的场所，这种积累后来会成为介入更大政治网络的资源。
\end{event}

\begin{event}{西周中晚期}{三门峡虢国墓地以车马、礼器与外来物资呈现要道节点}{考古遗存可靠；支系世系和具体职能不全确定}{三门峡及王畿相连的交通路线}\eventanchor{WZ-EARLY-GUO-CEMETERY}{西周·虢国墓地与要道网络}
三门峡虢国墓地中的车马、礼器和外来物资，使“封国”不再只是东向展开的概念。王族支系据守的要道，需要以可见的仪式和资源来显示其地位，也需要把王畿与多条路线上的人、物和消息接起来。墓地因此提示网络有其空间厚度：地方中心既可接受王室关系，也能汇聚自身的交通优势。

考古材料并未把每一支虢的谱系和职责完整写在地上。丰富的随葬品不等于每一次物资流动都由王命调度，外来物品也不能自动说明固定贸易制度。它能证明的是连接和资源集中，不能替代对具体路线、职掌和时期差异的谨慎判断。
\end{event}

关于懿王“烹齐哀公”的故事，\eventanchor{WZ-MID-QI-AIGONG-TRADITION}{西周·齐哀公受惩的传统}保留在\sourcebook{史记·齐太公世家}及相关编年叙述中。它以严厉惩罚解释封国继承的紧张，因而适合用来观察后世史家如何想象王室与地方的关系；但记载孤立，人物动机和事件本身均有争议，不能把它写成一桩已由同时代材料证实的处决。即使故事不可坐实，它仍提醒我们，封国网络从来不是没有继承冲突和王室裁断风险的静态结构。

\begin{event}{西周末至东周初}{琉璃河的地方中心持续或重组}{遗存可以比较；不能把变化直接归因于镐京陷落}{燕地的城址、墓地与青铜器网络}\eventanchor{WZ-LATE-LIULIHE-CONTINUITY}{西周·琉璃河的跨期延续}
西周王室的危机不等于每一个封国中心同时消失。琉璃河晚期材料和燕国青铜器所呈现的，是地方中心可能持续、调整或重新组合的轮廓。城址、墓地与地方资源使北方节点具有自己的时间尺度；它们不必随镐京政治的每一次震荡同步中断。

这种跨期材料也有严格限度。遗存的连续性不能直连为某一具体政治选择，更不能用来证明燕地在前771年之后采取了何种立场。可见的是中心并未必然空白；不可见的是每一段制度和人群关系怎样变化。将两者区分，才能避免把“网络解体”误写成所有地方能力一并消失。
\end{event}

\begin{event}{西周末至东周初}{天马—曲村墓地的跨期仪式延续}{连续性可靠；其对具体政治事件的影响无法逐墓判断}{晋地及其通向东周政治的资源基础}\eventanchor{WZ-LATE-TIANMA-QUCUN}{西周·天马—曲村的跨期仪式}
天马—曲村跨期墓地所见的祭祀连续，为晋地提供了另一条不同于王室编年的时间线。一个拥有土地、交通与家族仪式资源的宗亲中心，即使在西周终局中面对新的政治条件，仍可能保有组织和动员的基础。此后晋能进入东周事务的条件，不能只从春秋时代的结果倒推，也应在这种长期积累中寻找。

不过，墓地的连续只能证明仪式和精英记忆的延续，不能逐座墓解释晋国后来每一次联盟和战争。它为“具有基础”提供证据，不能把基础直接等同于结局。正是在这条界线上，考古材料与后来的国别叙事才可以互相校正，而不互相取代。
\end{event}
